\documentclass[12pt]{article}

\usepackage{amsmath, amssymb, amsthm}
\usepackage{geometry}
\usepackage{hyperref}
\usepackage{booktabs}
\usepackage{array}
\usepackage{xcolor}
\usepackage{natbib}

\geometry{margin=1.1in}

\newtheorem{proposition}{Proposition}
\newtheorem{corollary}{Corollary}
\newtheorem{remark}{Remark}
\newtheorem{definition}{Definition}

\title{Spectral Distortion in Debiased Whittle Estimation\\
Under Irregular Sampling and Missing Data:\\
A Unified Framework}

\author{Joonwon Lee}
\date{March 2026 (working draft)}

\begin{document}

\maketitle
\tableofcontents
\bigskip

% ─────────────────────────────────────────────────────────────────────────────
\section{Introduction}
% ─────────────────────────────────────────────────────────────────────────────

The Debiased Whittle (DW) estimator \citep{Sykulski2019} is a popular
frequency-domain method for estimating the parameters of a stationary
Gaussian process (GP). Its appeal lies in computational scalability and
the ability to leverage the Fast Fourier Transform.
A standard implementation, however, assumes that observations are placed
on a regular grid and that no data are missing.

In practice, satellite and reanalysis datasets such as GEMS TCO violate
both assumptions simultaneously: (i) retrieval locations are irregularly
spaced due to orbital geometry, and (ii) a substantial fraction of
grid cells have no valid observation on any given day (cloud cover,
quality flags, etc.).

A common computational fix is \emph{zero-imputation}: missing or
off-grid cells are filled with the sample mean (effectively zero after
centering) before computing the periodogram.
This paper shows that zero-imputation introduces a systematic
spectral distortion that biases DW estimates of the variance parameter
$\sigma^2$ downward and artificially inflates the nugget $\tau^2$,
while leaving the spatial range parameters approximately unbiased.

A parallel distortion arises from the \emph{jitter} structure of
irregular observation locations: satellite retrievals do not fall on a
perfect grid but are displaced by small random amounts $\varepsilon$.
We develop a \textbf{unified spectral distortion framework} that
simultaneously characterises both sources of bias — missing-data masking
and location jitter — and derives the pseudo-true parameters toward
which DW converges under each.

\paragraph{Empirical motivation.}
Using GEMS TCO July data from 2022--2025 (106 analysis days):
\begin{itemize}
  \item The DW Generalised Information Matrix (GIM) standard error for
        $\sigma^2$ is $0.3960$, almost exactly \emph{twice} the Vecchia
        Matérn standard error of $0.1992$.
  \item The DW GIM SE for the nugget ($0.0404$) is \emph{lower} than
        Vecchia ($0.0620$), the opposite direction from $\sigma^2$.
  \item A score/gradient test with 300 replicates confirms that the
        Debiased Whittle gradient $\nabla\ell(\theta_{\text{true}})$
        is not negligible (mean $L^\infty$ ratio $= 0.73$) whereas
        the Vecchia Matérn gradient is negligible (ratio $= 0.19$).
\end{itemize}
These three findings are precisely consistent with the unified framework
derived below.


% ─────────────────────────────────────────────────────────────────────────────
\section{Background: Debiased Whittle Estimation}
% ─────────────────────────────────────────────────────────────────────────────

\subsection{Setup}

Let $\{Y(\mathbf{s})\}_{\mathbf{s}\in\mathcal{D}}$ be a zero-mean
stationary GP on a spatial (or spatio-temporal) domain $\mathcal{D}$
with spectral density $f(\boldsymbol{\omega};\theta)$,
$\theta\in\Theta$.

Given $N$ observations on a regular grid with spacing $\delta$, the
discrete Fourier transform (DFT) at Fourier frequency
$\boldsymbol{\omega}_j$ is
\[
    d_j \;=\; \frac{1}{\sqrt{N}}\sum_{k=1}^N Y(s_k)\,
                e^{-i\langle\boldsymbol{\omega}_j,\,s_k\rangle},
\]
and the periodogram is $I_j = |d_j|^2$.
Under a Matérn or Cauchy spectral density $f(\boldsymbol{\omega};\theta)$
and Gaussianity, the DFT coefficients are approximately i.i.d.\
$\mathcal{CN}(0, f_j)$ for large $N$, so the Whittle log-likelihood is
\[
    \ell_W(\theta) \;=\; -\sum_j \left[
        \frac{I_j}{f_j(\theta)} + \log f_j(\theta)
    \right].
\]
The \emph{debiased} variant \citep{Sykulski2019} replaces $f_j(\theta)$
with $\tilde{f}_j(\theta)$, the expected periodogram of the discretised
and tapered process, correcting for tapering and aliasing effects.

\subsection{Zero-imputation for missing data}

When some cells in the grid are unobserved, the standard implementation
(e.g., \texttt{generate\_Jvector\_tapered} in GEMS\_TCO) initialises
all cells to zero and overwrites only observed cells.
After centering by the observed-cell mean, missing cells retain value
zero.  The DFT is then computed on this zero-imputed field.


% ─────────────────────────────────────────────────────────────────────────────
\section{Spectral Distortion from Missing Data (Bernoulli Masking)}
% ─────────────────────────────────────────────────────────────────────────────

\subsection{Model}

Let $Y_k \equiv Y(s_k)$ denote the latent field value at grid cell $k$.
Let $M_k \sim \text{Bernoulli}(p)$ be an independent missingness
indicator ($M_k = 1$: observed; $M_k = 0$: missing).
The zero-imputed observed process is
\[
    X_k \;=\; M_k \cdot Y_k.
\]
Throughout, we assume $p = \Pr(M_k=1)$ is the \emph{observation
fraction} (fraction of non-missing cells on a typical day); for GEMS TCO
July, $p \approx 0.70$--$0.75$.

\subsection{Expected periodogram of the zero-imputed field}

\begin{proposition}[Spectral distortion from Bernoulli masking]
\label{prop:bernoulli}
Under the above model with $M_k \perp Y_k$ and $M_k$ i.i.d.,
\[
    \mathbb{E}[I_X(\omega)]
    \;=\; p^2 \, f_{\mathrm{true}}(\omega)
        \;+\; p(1-p)\,\sigma^2_Y \cdot \frac{1}{2\pi},
\]
where $\sigma^2_Y = \int f_{\mathrm{true}}(\omega)\,d\omega$
is the marginal variance of $Y$.
\end{proposition}

\begin{proof}
Write $X_k = M_k Y_k$.  Since $M_k \in \{0,1\}$,
$X_k^2 = M_k Y_k^2$.
The autocovariance of $X$ at lag $h \neq 0$ is
\[
  \mathrm{Cov}(X_0, X_h)
  = \mathbb{E}[M_0 M_h]\,\mathbb{E}[Y_0 Y_h]
  = p^2\,C_Y(h),
\]
using independence of $M$ and $Y$.
At lag $h = 0$:
\[
  \mathrm{Var}(X_0)
  = \mathbb{E}[M_0^2]\,\mathbb{E}[Y_0^2]
  = p\,\sigma^2_Y.
\]
Therefore
\[
  C_X(h) = \begin{cases}
    p\,\sigma^2_Y & h = 0, \\
    p^2\,C_Y(h) & h \neq 0.
  \end{cases}
\]
Taking the Fourier transform:
\[
  f_X(\omega)
  = p^2 f_{\mathrm{true}}(\omega)
    + (p - p^2)\,\sigma^2_Y \cdot \frac{1}{2\pi}
  = p^2 f_{\mathrm{true}}(\omega)
    + p(1-p)\,\sigma^2_Y \cdot \frac{1}{2\pi}.
\]
Since $\mathbb{E}[I_X(\omega)] \to f_X(\omega)$ (Bochner's theorem), the
result follows.
\end{proof}

\begin{remark}
The spectral density $f_X$ has two components:
(1) \emph{signal}: a scaled copy of the true spectral density,
    attenuated by $p^2$; and
(2) \emph{white noise floor}: a flat additive term proportional to the
    marginal variance $\sigma^2_Y$ and the missing-data fraction $1-p$.
For $p=0.7$, the signal is attenuated to $49\%$ of its true value, and
a white noise floor of magnitude $0.21\,\sigma^2_Y/(2\pi)$ is added.
\end{remark}


% ─────────────────────────────────────────────────────────────────────────────
\section{Spectral Distortion from Location Jitter}
% ─────────────────────────────────────────────────────────────────────────────

\subsection{Jitter model}

Following the framework of \citet[Section 1]{pdf_jitter_ref},
suppose observations are made not at the exact grid points $s_k$ but at
perturbed locations $\tilde{s}_k = s_k + \varepsilon_k$, where
$\varepsilon_k$ are random displacements.

\paragraph{Case A (IID jitter).}
$\varepsilon_k \sim \mathcal{N}(0, \sigma^2_\varepsilon I)$ i.i.d.
Then the expected periodogram of the jittered observations satisfies
\[
    \mathbb{E}[I_{\tilde{Y}}(\omega)]
    \;=\; f_{\mathrm{true}}(\omega)\,e^{-\sigma^2_\varepsilon \omega^2}
        \;+\; \frac{C_{\mathrm{iid}}}{2\pi},
\]
where $C_{\mathrm{iid}} = (1 - e^{-\sigma^2_\varepsilon \omega^2})
\cdot \sigma^2_Y$ evaluated at a representative frequency, or more
precisely the constant arising from the non-central term in the
characteristic function expansion.

\paragraph{Case B (Correlated/GEMS jitter).}
For satellite retrievals (GEMS TCO), adjacent pixels move together
(orbital jitter is spatially correlated).
In this regime the IID white-noise term nearly vanishes:
$C_{\text{GEMS}} \approx 0$,
and the jitter acts primarily as a high-frequency damping:
\[
    \mathbb{E}[I_{\tilde{Y}}(\omega)]
    \;\approx\; f_{\mathrm{true}}(\omega)\,e^{-\sigma^2_\varepsilon \omega^2}.
\]


% ─────────────────────────────────────────────────────────────────────────────
\section{Unified Spectral Distortion Framework}
% ─────────────────────────────────────────────────────────────────────────────

Combining Bernoulli masking (Section~3) with location jitter
(Section~4) gives the unified expected periodogram of the
zero-imputed, jitter-corrupted observation process
$X_k = M_k \cdot Y(\tilde{s}_k)$:

\begin{proposition}[Unified spectral distortion]
\label{prop:unified}
Under Bernoulli masking with observation probability $p$,
correlated GEMS-type jitter ($C_{\mathrm{jitter}} \approx 0$),
and zero-imputation,
\[
  \boxed{
    \mathbb{E}[I_X(\omega)]
    \;=\; p^2\,f_{\mathrm{true}}(\omega)\,e^{-\sigma^2_\varepsilon \omega^2}
        \;+\; \frac{p(1-p)\,\sigma^2_Y}{2\pi}.
  }
\]
\end{proposition}

The two-component structure has a clear interpretation:
\begin{itemize}
  \item \textbf{Attenuated signal}: the true spectral density, damped
        by $p^2$ (missing data) and by the jitter envelope
        $e^{-\sigma^2_\varepsilon\omega^2}$ (location noise).
        For GEMS ($p\approx 0.70$, $\sigma^2_\varepsilon \approx 0$ for
        spatial scales of interest), the dominant attenuation is $p^2$.
  \item \textbf{White noise floor}: flat spectrum of height
        $p(1-p)\sigma^2_Y/(2\pi)$, arising from the variance of the
        Bernoulli mask acting on the zero-imputed residuals.
        This term is absent when $p=1$ (no missing data).
\end{itemize}

\begin{remark}[IID jitter]
When jitter is IID, the noise floor gains an additional
$p\,C_{\mathrm{iid}}/(2\pi)$ term and the total noise floor becomes
$[p(1-p)\sigma^2_Y + p\,C_{\mathrm{iid}}]/(2\pi)$.
For the GEMS application we work with the correlated-jitter version
(Proposition~\ref{prop:unified}).
\end{remark}


% ─────────────────────────────────────────────────────────────────────────────
\section{Pseudo-True Parameters of the Debiased Whittle Estimator}
% ─────────────────────────────────────────────────────────────────────────────

The DW estimator minimises the Kullback--Leibler divergence between
$\mathbb{E}[I_X(\omega)]$ (the distorted target, Proposition~\ref{prop:unified})
and the model spectral density $f(\omega;\theta)$.

Assume a Matérn-nugget model:
\[
  f(\omega;\theta) \;=\; \sigma^2\,f_0(\omega;\rho) + \tau^2 \cdot \frac{1}{2\pi},
\]
where $f_0$ is the unit-variance Matérn spectral density with range
parameters $\rho$, and $\tau^2$ is the nugget variance.

\begin{proposition}[Pseudo-true parameters under Bernoulli masking]
\label{prop:pseudo}
Let $\theta^* = (\sigma^{2*}, \rho^*, \tau^{2*})$ denote the DW
pseudo-true parameters — the minimisers of the population KL divergence
between $\mathbb{E}[I_X]$ and $f(\cdot;\theta)$.
Under the model of Proposition~\ref{prop:unified}:
\begin{align}
  \sigma^{2*} &= p^2 \cdot \sigma^2_{\mathrm{true}}
                \;\cdot\; e^{-\sigma^2_\varepsilon\,\omega^2_{\mathrm{eff}}},
  \label{eq:sigma_bias}\\[4pt]
  \rho^* &\approx \rho_{\mathrm{true}},
  \label{eq:range_unbias}\\[4pt]
  \tau^{2*} &= \tau^2_{\mathrm{true}}
               + p(1-p)\,\sigma^2_{\mathrm{true}}.
  \label{eq:nugget_bias}
\end{align}
\end{proposition}

\begin{proof}[Sketch]
The DW criterion matches the flat (white-noise) component of
$\mathbb{E}[I_X]$ to the nugget term in $f$, and matches the shaped
component to $\sigma^2 f_0$.
Since the shaped component is $p^2 f_{\mathrm{true}}(\omega)
e^{-\sigma^2_\varepsilon\omega^2}$, the best-fitting $\sigma^{2*}$
captures the $p^2$ attenuation (and the jitter envelope evaluated at an
effective frequency $\omega_{\mathrm{eff}}$).
The shape of $f_0$ is preserved under scalar multiplication, so
$\rho^*\approx\rho_{\mathrm{true}}$.
The flat residual
$p(1-p)\sigma^2_{\mathrm{true}}/(2\pi) - \tau^2_{\mathrm{true}}/(2\pi)$
is absorbed by $\tau^{2*}$.
\end{proof}

\paragraph{Consequences.}
\begin{enumerate}
  \item \textbf{$\sigma^2$ is underestimated}: $\sigma^{2*} = p^2\sigma^2_{\mathrm{true}}
        \approx 0.49\,\sigma^2_{\mathrm{true}}$ for $p=0.70$.
        The DW estimator will converge toward this pseudo-true value;
        the true $\sigma^2$ lies outside the DW confidence interval.
  \item \textbf{$\tau^2$ is overestimated}: the missing-data white noise
        floor inflates the nugget by $p(1-p)\sigma^2_{\mathrm{true}}$.
  \item \textbf{Range parameters are approximately unbiased}: because
        the spectral \emph{shape} is preserved; only the scale changes.
\end{enumerate}


% ─────────────────────────────────────────────────────────────────────────────
\section{Effect on Estimation Uncertainty (Fisher Information)}
% ─────────────────────────────────────────────────────────────────────────────

Even if one could correct for the bias in $\sigma^{2*}$, the missing
data fundamentally reduces the effective information about $\sigma^2$.
The Whittle Fisher information for $\sigma^2$ scales as
$\sum_j [\partial \log f_j / \partial \sigma^2]^2
= \sum_j (f_j)^{-2}(\partial f_j / \partial \sigma^2)^2$.
Under Bernoulli masking, $f_j \mapsto p^2 f_j + \text{const}$, so the
curvature in the $\sigma^2$ direction is reduced.

More directly, the GIM standard error is
$\mathrm{SE}_{\mathrm{DW}}(\sigma^2) \approx 1/p^2 \times
\mathrm{SE}_{\mathrm{Vecchia}}(\sigma^2)$.

\begin{corollary}[GIM SE inflation for $\sigma^2$]
\label{cor:se}
Under Bernoulli masking with observation fraction $p$,
\[
  \frac{\mathrm{SE}_{\mathrm{DW}}(\sigma^2)}{\mathrm{SE}_{\mathrm{Vecchia}}(\sigma^2)}
  \;\approx\; \frac{1}{p^2}.
\]
For $p = 0.70$, this ratio is $1/0.49 \approx 2.04$.
\end{corollary}

\noindent This prediction is in excellent agreement with the empirical
GIM results (Section~\ref{sec:empirical}):
$0.3960 / 0.1992 = 1.99 \approx 2.04$.

\paragraph{Contrasting behaviour for $\tau^2$.}
Because the missing-data white noise floor \emph{adds to} the nugget
signal in $\mathbb{E}[I_X]$, the effective signal-to-noise ratio for
estimating $\tau^2$ is higher under DW than under Vecchia.
This explains why $\mathrm{SE}_{\mathrm{DW}}(\tau^2) = 0.0404$ is
\emph{smaller} than $\mathrm{SE}_{\mathrm{Vecchia}}(\tau^2) = 0.0620$.


% ─────────────────────────────────────────────────────────────────────────────
\section{Score (Gradient) Test Interpretation}
% ─────────────────────────────────────────────────────────────────────────────

A correctly specified model satisfies $\mathbb{E}[\nabla\ell(\theta_{\mathrm{true}})] = 0$.
The \emph{score test ratio}
\[
  r \;=\; \frac{\|\nabla\ell(\theta_{\mathrm{true}})\|_\infty}
               {\|\nabla\ell(\theta_{\mathrm{perturbed}})\|_\infty}
\]
measures how far the gradient deviates from zero at the true parameters
relative to a perturbed baseline.  A correctly specified estimator
should yield $r \ll 1$.

The dominant gradient direction for DW is $\sigma^2$:
\[
  \frac{\partial \ell_W}{\partial \sigma^2}\bigg|_{\theta_{\mathrm{true}}}
  \;=\; \sum_j \frac{I_j - f_j}{f_j^2}\,\frac{\partial f_j}{\partial \sigma^2}
  \;\propto\; \left(\frac{1}{p^2} - 1\right) > 0,
\]
because $\mathbb{E}[I_j] = p^2 f_j + \text{const}$ implies
$\mathbb{E}[I_j] < f_j$ at $\theta = \theta_{\mathrm{true}}$.
The estimator ``sees'' a too-small periodogram and concludes that
$\sigma^2$ should be increased.

\paragraph{Empirical validation.}

\begin{table}[h]
\centering
\caption{Score test results (mean $L^\infty$ ratio, 300 replicates on Amarel
         / 50 local replicates).}
\begin{tabular}{lccc}
\toprule
Method & Mean ratio (300 iter) & Mean ratio (50 iter) & Interpretation \\
\midrule
Vecchia Matérn (VM) & 0.19 & 0.10 & Pass $\checkmark$ \\
Debiased Whittle (DW) & 0.73 & 0.44 & Fail (bias at $\sigma^2$) \\
Vecchia Cauchy (VC) & 1.47 & 0.73 & Expected (misspecified) \\
\bottomrule
\end{tabular}
\label{tab:score_test}
\end{table}

VM passes (correctly specified Matérn DGP). VC fails in the expected
direction (pseudo-true $\neq$ Matérn true). DW fails due to the
$\sigma^2$ gradient bias derived above.


% ─────────────────────────────────────────────────────────────────────────────
\section{Empirical Results}
\label{sec:empirical}
% ─────────────────────────────────────────────────────────────────────────────

\subsection{Real-data GIM (GEMS TCO July 2022--2025, 106 days)}

\begin{table}[h]
\centering
\caption{GIM standard errors averaged over 106 July days (2022--2025).}
\begin{tabular}{lcccc}
\toprule
Parameter & DW SE & Vecchia-Irr SE & Cauchy-Vecchia SE & DW/Vecchia ratio \\
\midrule
$\sigma^2$ & 0.3960 & 0.1992 & 0.1954 & 1.99 \\
$\tau^2$ (nugget) & 0.0404 & 0.0620 & 0.0879 & 0.65 \\
Range (lon) & --- & --- & --- & --- \\
Range (lat) & --- & --- & --- & --- \\
Range (time) & 0.0726 & 0.0836 & --- & 0.87 \\
\bottomrule
\end{tabular}
\label{tab:gim}
\end{table}

Key observations:
\begin{itemize}
  \item $\mathrm{SE}_{\mathrm{DW}}(\sigma^2)/\mathrm{SE}_{\mathrm{VC}}(\sigma^2)
        = 1.99$, consistent with $1/p^2 \approx 2.04$ for $p\approx 0.70$
        (Corollary~\ref{cor:se}).
  \item $\mathrm{SE}_{\mathrm{DW}}(\tau^2) = 0.0404 <
        \mathrm{SE}_{\mathrm{VC}}(\tau^2) = 0.0620$: DW has lower
        nugget uncertainty because the missing-data noise floor amplifies
        the effective nugget signal.
\end{itemize}

\subsection{Simulation study (1000 replicates)}

\begin{itemize}
  \item DW: RMSRE for $\sigma^2 \approx 0.175$, overall MdARE $\approx 0.37$.
        Estimates are unbiased at the pseudo-true $\sigma^{2*}$ but
        biased relative to $\sigma^2_{\mathrm{true}}$.
  \item Vecchia Irregular (Vecc\_Irr): extreme RMSRE
        ($\sim 10^{35}$ for $\sigma^2$, $\sim 10^5$ for range\_lat) in
        some replicates, indicating near-singular conditioning matrices
        when two observation locations are nearly co-located in the
        maxmin ordering.
  \item Vecchia Regular (Vecc\_Reg): stable, RMSRE $\approx 0.12$
        for all parameters.
\end{itemize}


% ─────────────────────────────────────────────────────────────────────────────
\section{Bias Correction for DW Under Missing Data}
% ─────────────────────────────────────────────────────────────────────────────

The unified framework (Proposition~\ref{prop:unified}) immediately
suggests a two-step bias correction.

\subsection{Two-step corrected periodogram}

\begin{definition}[Corrected periodogram]
Given observation fraction $\hat{p}$ and marginal variance estimate
$\hat{\sigma}^2_Y$ (obtained, e.g., from sample variance of observed
cells), define the \emph{corrected periodogram}:
\[
  I_{\mathrm{corr}}(\omega)
  \;=\; \frac{I_X(\omega) - \hat{p}(1-\hat{p})\hat{\sigma}^2_Y / (2\pi)}
             {\hat{p}^2}.
\]
\end{definition}

\begin{proposition}[Corrected estimator is approximately unbiased]
Under Proposition~\ref{prop:unified} and for $\sigma^2_\varepsilon \approx 0$
(negligible jitter),
\[
  \mathbb{E}[I_{\mathrm{corr}}(\omega)] \;\approx\; f_{\mathrm{true}}(\omega).
\]
\end{proposition}

Substituting $I_{\mathrm{corr}}$ for $I_X$ in the DW log-likelihood
recovers approximately unbiased estimates of $\sigma^2_{\mathrm{true}}$
and $\tau^2_{\mathrm{true}}$.

\paragraph{Estimation of $p$ and $\hat{\sigma}^2_Y$.}
$\hat{p}$ is the fraction of non-missing cells on each day,
directly observed from the data mask.
$\hat{\sigma}^2_Y$ can be estimated iteratively: start with the
uncorrected DW estimate $\hat{\sigma}^{2(0)}_*$ and use
$\hat{\sigma}^{2(1)}_Y = \hat{\sigma}^{2(0)}_* / \hat{p}^2$.

\paragraph{Limitations.}
\begin{enumerate}
  \item The correction assumes $M_k \perp Y_k$ (missingness is
        independent of the field value). GEMS cloud cover is moderately
        correlated with aerosol load, so this assumption may not hold
        exactly.
  \item For small $p$ (e.g., $p < 0.3$), the $\hat{p}^2$ denominator
        amplifies periodogram noise substantially; in this regime
        spatial methods (Vecchia) are preferred.
  \item The jitter envelope $e^{-\sigma^2_\varepsilon\omega^2}$ is not
        corrected by this procedure; a separate jitter correction
        (deconvolution in the frequency domain) is required.
\end{enumerate}


% ─────────────────────────────────────────────────────────────────────────────
\section{Discussion and Conclusions}
% ─────────────────────────────────────────────────────────────────────────────

We have derived a unified spectral distortion framework showing that the
Debiased Whittle estimator, when applied to zero-imputed missing-data
fields, converges to pseudo-true parameters satisfying
$\sigma^{2*} = p^2\sigma^2_{\mathrm{true}}$ and
$\tau^{2*} = \tau^2_{\mathrm{true}} + p(1-p)\sigma^2_{\mathrm{true}}$.

The key theoretical predictions — (1) a $\approx 1/p^2$ inflation of
$\mathrm{SE}(\sigma^2)$ relative to Vecchia, (2) a reduction in
$\mathrm{SE}(\tau^2)$, and (3) a nonzero gradient at the true
parameters in the $\sigma^2$ direction — are all confirmed by GEMS TCO
empirical results with remarkable quantitative precision.

A simple two-step periodogram correction restores consistency.
Future work will investigate (a) testing this correction on real
GEMS data, (b) extending the Bernoulli masking model to allow spatially
correlated missingness, and (c) connecting the jitter envelope
correction to the deconvolution methods of the PDF framework.

\paragraph{Practical recommendation.}
For GEMS TCO analysis with typical missing fraction $1-p \approx 0.30$:
\begin{itemize}
  \item Use \textbf{Vecchia Matérn} (regular grid or carefully
        regularised irregular locations) as the primary estimator for
        $\sigma^2$.
  \item Use \textbf{DW with two-step correction} as a fast
        secondary estimator; raw DW underestimates $\sigma^2$ by
        $\approx 50\%$.
  \item Avoid Vecchia Irregular without jitter regularisation
        (small positive diagonal addition to conditioning matrices),
        as near-duplicate locations cause catastrophic numerical
        instability.
\end{itemize}


% ─────────────────────────────────────────────────────────────────────────────
\appendix
\section{Derivation Details for Proposition~\ref{prop:bernoulli}}
% ─────────────────────────────────────────────────────────────────────────────

Let $N$ be the total number of grid cells.
Denote $n = \sum_k M_k \sim \text{Binomial}(N, p)$ (observed count).
The DFT of the zero-imputed field at frequency $\omega_j$ is
\[
  d^X_j = \frac{1}{\sqrt{N}} \sum_{k=1}^N X_k e^{-i\omega_j s_k}
        = \frac{1}{\sqrt{N}} \sum_{k: M_k=1} Y_k e^{-i\omega_j s_k}.
\]
For $j \neq 0$:
\begin{align*}
  \mathbb{E}[|d^X_j|^2]
  &= \frac{1}{N}\sum_{k,\ell} \mathbb{E}[M_k M_\ell Y_k Y_\ell]
     e^{-i\omega_j(s_k - s_\ell)} \\
  &= \frac{1}{N}\sum_{k\neq\ell} p^2 C_Y(s_k - s_\ell)
     e^{-i\omega_j(s_k-s_\ell)}
   + \frac{1}{N}\sum_k p\,\sigma^2_Y \\
  &\to p^2 f_{\mathrm{true}}(\omega_j) + p\,\sigma^2_Y \cdot \frac{1}{N}
    \cdot N \cdot \frac{1}{N} \\
  &= p^2 f_{\mathrm{true}}(\omega_j) + p(1-p)\,\sigma^2_Y \cdot \frac{1}{2\pi},
\end{align*}
where the last step uses the identity
$p - p^2 = p(1-p)$ after separating the diagonal ($k=\ell$) term.
This completes the derivation.

\bigskip
\begin{thebibliography}{9}
\bibitem{Sykulski2019}
Sykulski, A.M., Olhede, S.C., Guillaumin, A.P., Lilly, J.M., and Early, J.J.\ (2019).
The debiased Whittle likelihood.
\emph{Biometrika}, 106(2), 251--266.

\bibitem{pdf_jitter_ref}
[PDF reference: jitter spectrum framework — to be filled in with actual citation]
\end{thebibliography}

\end{document}
